\section{Problem Formulation}
\label{sec:problem-formulation}

\subsection{Conventions and Notation}
\label{sec:conventions-notation}

We adopt the following conventions throughout the paper.
\begin{itemize}
  \item \textbf{Indexing.} $t\in\mathcal{T}$ denotes a task (panoramic image). $w\in\mathcal{W}$ denotes a worker/annotator.
  \item \textbf{Estimates.} A hat (e.g., $\hat{y}_t$) denotes a consensus estimate.
  \item \textbf{Missingness.} For meta-label entries, NA denotes a system/format missing value (e.g., field not transmitted). NA is \emph{never} silently dropped; it is reported as an auditable rate.
  \item \textbf{Pre-registered control knobs.} Protocol locking is staged: Pilot fixes immutable pool boundaries, sampling schemas, candidate grids, and split rules; PreScreen fixes admission and weighting thresholds; the first Calibration round fixes scene-routing thresholds that depend on realized coverage.
\end{itemize}

Table~\ref{tab:notation} summarizes the main symbols.
\small
\renewcommand{\arraystretch}{1.12}
\begin{longtable}{@{}p{0.22\linewidth}p{0.72\linewidth}@{}}
\caption{Key notation used in the problem formulation and system design.}\label{tab:notation}\\
	oprule
Symbol & Meaning \\
\midrule
\endfirsthead
\caption[]{Key notation used in the problem formulation and system design (continued).}\\
	oprule
Symbol & Meaning \\
\midrule
\endhead
$\mathcal{T}$ & Set of annotation tasks (panoramic images), $|\mathcal{T}|=N$. \\
$\mathcal{P}$ & Model proposal set, where $P_t\in\mathcal{P}$ is the initialization for task $t$. \\
$\mathcal{W}$ & Worker pool (Label Studio), $|\mathcal{W}|=M$. \\
$\pi$ & Task distribution strategy mapping tasks to assigned workers (budget-aware, stress-mode capable). \\
$B$ & Global assignment budget (e.g., total assignments across tasks). \\
$e_{w,t}$ & Active annotation effort for assignment $(w,t)$ (instrumented engagement time). \\
$\mathbf{m}_{w,t}$ & Meta-label vector (scope, difficulty tags, model-issue tags) for assignment $(w,t)$. \\
$r_w,\ r_w^{(s)}$ & Global and stratum-specific worker reliability estimates. \\
$\mathcal{C}$ & Calibration \\\& audit procedures (reliability estimation, counter-example detection). \\
$\mathcal{A}$ & Aggregation operator producing $\hat{y}_t$ from submissions. \\
$z_{t,w,k}$ & Binary tag indicator for task $t$, worker $w$, tag $k$; value in $\{0,1,\text{NA}\}$. \\
$c_{t,k}$ & Consensus decision for tag $k$ on task $t$ (binary; NA if undefined). \\
$\tau_c$ & Consensus threshold for per-tag majority (pre-registered and fixed before Main). \\
$\mathrm{IAA}_t$ & Intra-task agreement proxy (median pairwise IoU among submissions on task $t$). \\
$C_t$ & Task-level consensus representative submission (max-median agreement). \\
$C_t^{(-w)}$ & Leave-one-out consensus representative for evaluating worker $w$ on task $t$. \\
$\mathrm{IoU}_{\mathrm{LOO}}(t,w)$ & IoU between worker $w$'s submission and $C_t^{(-w)}$. \\
$\omega_w$ & Pre-computed reliability weight for worker $w$ (independent of the evaluated task). \\
$d_t,\ g_t$ & Pre-annotation risk proxies (distribution shift and structural failure signals). \\
$k,\ k_{\max}$ & Redundancy per task and its pre-registered upper bound. \\
$k_{\min}$ & Minimum redundancy required for calibration tasks in reliability estimation (pre-registered). \\
$k_0(t)$ & Initial redundancy for task $t$ selected by the routing policy (risk-stratified). \\
$\mathrm{LCB}(r_w)$ & Lower confidence bound of worker reliability used for conservative routing. \\
$R_w^{\mathrm{tier}}$ & Discrete worker risk tier used in profiling and routing audit ($R0$ stable, $R1$ trust-vulnerable, $R2$ condition-fragile, $R3$ unusable/noisy). \\
$T_w,\ C_w,\ G_w$ & Core risk signatures: blind-trust risk, condition fragility, and gate-failure tendency. \\
$\tau_{\mathrm{low}},\ \tau_{\mathrm{high}}$ & Routing thresholds on $\mathrm{LCB}(r_w)$ for default vs. stress mode. \\
$\tau_{\mathrm{stress}}$ & Stress-mode trigger threshold (pre-registered). \\
\bottomrule
\end{longtable}

\subsection{System Model}

We model the semi-automatic panoramic layout annotation process as a system that transforms a set of inputs---raw panoramic images $\mathcal{T}=\{t_1,\dots,t_N\}$ and initial model proposals $\mathcal{P}=\{P_1,\dots,P_N\}$---into a set of high-quality consensus labels $\{\hat{y}_t\}$ through a labor pool $\mathcal{W}=\{w_1,\dots,w_M\}$ with heterogeneous capabilities. 

The internal logic of this system is governed by a pre-registered protocol consisting of three components:
\begin{enumerate}
    \item \textbf{Annotation Task Distribution Strategy ($\pi$)}: A mapping $\pi(t) \subseteq \mathcal{W}$ that determines which annotators are assigned to task $t$, subject to a global assignment budget $B$:
    \begin{equation}
      \sum_{t\in\mathcal{T}}|\pi(t)| \le B.
    \end{equation}
    \item \textbf{Calibration and Audit Procedures ($\mathcal{C}$)}: A set of functions used to estimate worker reliability $r_w$ and detect OOD risks before or during the annotation process.
    \item \textbf{Aggregation Operator ($\mathcal{A}$)}: A deterministic procedure that synthesizes individual submissions $(y_{w,t}, \mathbf{m}_{w,t})$ into the consensus output $\hat{y}_t$.
\end{enumerate}

For each assignment $(w,t)$, the system observes a sample annotation effort $e_{w,t}\in\mathbb{R}_{\ge0}$ (instrumented as active engagement time) and a structured meta-label vector
\begin{equation}\label{eq:meta}
  \mathbf{m}_{w,t}=(\texttt{scope},\;\texttt{difficulty},\;\texttt{model\_issue}).
\end{equation}

The central design objective is to optimize the protocol parameters $(\pi, \mathcal{C}, \mathcal{A})$ to maximize the aggregate annotation quality $R$ across the full data distribution, while satisfying the budget constraint and the strict requirement that all rules are fixed prior to data collection:
\begin{equation}\label{eq:obj}
  \max_{\pi,\,\mathcal{C},\,\mathcal{A}}\quad R(\pi,\mathcal{C},\mathcal{A})\quad\text{s.t.}\quad
  \sum_{t}|\pi(t)|\le B.
\end{equation}

All protocol decisions (i.e., the choice of $(\pi,\mathcal{C},\mathcal{A})$ and associated thresholds) follow a staged freeze schedule: Pilot fixes the protocol skeleton and non-overlap constraints, PreScreen fixes admission-related values, and only after the first Calibration round are scene-routing thresholds frozen. This prevents post-hoc freedom while avoiding premature locking of parameters that require calibration coverage evidence.

\subsection{Three Perspectives on System Performance}
\label{sec:dimensions}

The evaluation of \eqref{eq:obj} is decomposed into three perspectives that address the efficiency, reliability, and robustness of the proposed framework.

\begin{itemize}
  \item \textbf{Perspective 1 --- Efficiency improvement.}
  A fundamental requirement is to determine what constitutes a robust measure of data-annotation human cost
  under varying levels of task difficulty and model performance. We treat the sample annotation effort
  $e_{w,t}$ as a proxy for cognitive load. This perspective evaluates whether semi-automatic initialization
  reduces $e_{w,t}$ compared to manual baselines and how effort relates to latent image difficulty $d_t$.

  \item \textbf{Perspective 2 --- Annotation quality (robust consensus).}
  Annotation quality is quantified through the lens of robust consensus. We estimate worker reliability
  $r_w$ (and scene-specific variants $r_w^{(s)}$) to inform aggregation $\mathcal{A}$.
  Any reliability-based weighting must be validated on held-out calibration data to avoid circular evaluation
  (i.e., using the same signal to both construct and judge $\hat{y}_t$).
  A key challenge is to distinguish genuine errors from irreducible task ambiguity; we therefore define
  diagnostic counter-examples that are explicitly audited:
  \begin{itemize}
    \item \textbf{Type 1 (Consensus divergence)}: internal agreement ($\mathrm{IAA}_t$) remains low despite sufficient redundancy.
    \item \textbf{Type 2 (Effort--outcome inconsistency)}: the relationship between $e_{w,t}$ and geometric change magnitude is anomalous.
    \item \textbf{Type 3 (Geometric validity failure)}: structural violations such as self-intersecting polygons or topological collapse.
    \item \textbf{Type 4 (Process/format failure)}: missing required meta-fields in $\mathbf{m}_{w,t}$ (NA due to system/transport/export errors), or non-compliant meta-label records that violate UI-level conventions (e.g., empty selection when ``at least one'' is required, or mutually-exclusive tags co-selected such as \texttt{trivial}+\{other difficulty\} / \texttt{acceptable}+\{other failure modes\}).
  \end{itemize}

  \item \textbf{Perspective 3 --- Overall system efficiency.}
  This perspective evaluates the annotation task distribution strategy $\pi$ at scale under non-IID/OOD
  conditions. In practice, the system relies on pre-annotation risk proxies (available before labels are
  produced) to trigger a \emph{stress mode} that temporarily increases redundancy on suspected hard/OOD tasks.
  Using offline replay on tasks with $k$ real annotations, we trace the empirical trade-off between aggregate
  reliability $R$ and the average cost $\bar{k}=B/N$. Overall system efficiency is demonstrated if the protocol
  maintains $R$ across the full T/I/M distribution while minimizing redundant effort spent on easy, reliable
  tasks, without relying on silent pruning.
\end{itemize}

\subsection{Pre-registration and Auditability}

To ensure scientific rigor and to eliminate post-hoc freedom, the protocol follows a Pilot $\rightarrow$
PreScreen $\rightarrow$ Main procedure with staged locking rather than a single one-shot freeze. Pre-registration specifies:
\begin{enumerate}
  \item at Pilot: pool boundaries, non-overlap constraints, anchor/trap sampling schemas, candidate weighting grids, split rules, perturbation operators, and audit/logging fields;
  \item at PreScreen: the aggregation operator $\mathcal{A}$, admission thresholds, the candidate weighting family, and the cross-fitting rule used to select any weighting cap;
  \item for \texttt{PreScreen\_semi}, the source mix of misleading initializations (synthetic operators as the primary source, with a fixed small natural-failure supplement), together with non-overlap constraints between manual and semi pools and planned-vs.-realized quota disclosure;
  \item after the first Calibration round: scene-routing thresholds and activation rules that depend on realized coverage; and
  \item the diagnostics and disclosure tiers (T/I/M) together with the exact denominators used for reported rates.
\end{enumerate}

Evaluation uses offline replay with leave-one-task-out estimation and stratified bootstrap to
compute confidence intervals and sensitivity checks for all primary outcomes.

